\section{Conclusão}\label{sec:conclusao}

Este trabalho teve por objetivo desenvolver e avaliar o SIRAS, um sistema especialista baseado em
regras que interpreta a análise de solo, recomenda calagem e adubação para 61 culturas e grupos de
culturas segundo o Manual de Calagem e Adubação do RS e de SC e estima a aptidão edáfica da área.
O objetivo foi alcançado: as recomendações concordaram com o oráculo calculado à mão em todos os
25 casos de validação, com limite inferior do intervalo de confiança acima da meta de 80\%, o que
levou à aceitação da hipótese; a implementação da aptidão reproduziu a especificação congelada em
todos os 88 casos de conformidade; e o motor reproduziu as 32 classificações de fósforo e potássio
que o próprio Manual publica no Anexo~1.
Na avaliação de usabilidade, os oito participantes do público-alvo concluíram sozinhos um
cenário de recomendação, nenhum apontou problema além do estético, e todos disseram que usariam a
ferramenta no trabalho.

A principal contribuição é a integração, a partir de uma única análise de solo, de duas
perspectivas que as ferramentas disponíveis tratam em separado: a recomendação de correção e
adubação e a avaliação de aptidão, apresentada antes e depois da correção, de modo que o laudo
mostra não só quanto aplicar, mas o que a aplicação muda. A ela soma-se a rastreabilidade: cada
número do laudo indica a regra e a tabela do Manual que o produziram. A validação mostrou ainda
por que um oráculo independente importa: os casos acrescentados por lacuna de cobertura revelaram
dois defeitos que haviam passado por milhares de testes automatizados.

As limitações mais relevantes são três. A aptidão é exclusivamente edáfica, e seus critérios,
definidos pelo autor a partir do Manual e do sistema oficial de avaliação de aptidão, não foram
validados por especialistas independentes nem em campo. O oráculo das recomendações deriva do
mesmo Manual que o sistema implementa, o que limita a validação à fidelidade aos critérios
publicados. E o conjunto de casos é menor que o previsto e desigual entre os grupos de culturas.
A avaliação de usabilidade, com amostra pequena e por conveniência, indica ausência de obstáculos
graves, mas não permite generalizar a satisfação observada.

Como trabalhos futuros, destacam-se: a validação dos critérios de aptidão por um painel de
agrônomos e por dados de campo; a ampliação do conjunto de casos, sobretudo nos grupos com menos
casos, incluindo laudos reais com recomendação de profissional habilitado; a comparação sistemática
com ferramentas em uso, como o FertFácil, em resultados e em usabilidade; uma versão para celular
que funcione sem conexão, pedida por participante da avaliação; e a inclusão de culturas hoje fora do escopo e de manuais de outras regiões, que a
separação entre base de conhecimento e motor de inferência facilita.
